\documentclass[10pt,conference]{IEEEtran}
\IEEEoverridecommandlockouts

\usepackage[utf8]{inputenc}
\usepackage[T1]{fontenc}
\usepackage{amsmath,amssymb,amsfonts}
\usepackage{algorithmic}
\usepackage{algorithm}
\usepackage{graphicx}
\usepackage{booktabs}
\usepackage{multirow}
\usepackage{url}
\usepackage{cite}
\usepackage{xcolor}
\usepackage{tabularx}
\usepackage{array}

% Prevent overfull hboxes
\tolerance=1000
\emergencystretch=2em
\hfuzz=2pt
\sloppy

\begin{document}

\title{Quantum Kernel Circuit Design for Four-Class ECG Image
Classification: Encoding Range, Entanglement Topology,
and Statistical Validation}

\author{%
\IEEEauthorblockN{(Author 1 Name), (Author 2 Name), (Author 3 Name), (Author 4 Name)}
\IEEEauthorblockA{\textit{Dept.\ of Computer Science and Engineering}\\
St.\ Joseph Engineering College, Mangaluru, India\\
\{email1, email2, email3, email4\}@sjec.ac.in}
\and
\IEEEauthorblockN{(Guide Name)}
\IEEEauthorblockA{\textit{Dept.\ of Computer Science and Engineering}\\
St.\ Joseph Engineering College, Mangaluru, India}
}

\maketitle

\begin{abstract}
Quantum machine learning has been applied to automated ECG analysis, but
prior studies frequently report isolated accuracy figures without
parametric ablation, statistical validation, or confidence calibration.
This paper presents a hybrid classical-quantum pipeline for four-class
ECG image classification and systematically ablates five quantum circuit
variables. Features are extracted from a frozen ResNet50, compressed to
nine dimensions via TruncatedSVD, and classified with a nine-qubit
ZZFeatureMap fidelity kernel. On a 928-image clinical dataset, the
tuned Quantum Support Vector Classifier reaches 94.62\% accuracy. A
$2\!\times\!2$ factorial ablation shows that expanding the feature
encoding range to $[0,\pi]$ and closing the qubit chain into a circular
topology each improve accuracy, yielding a joint 4.84 percentage-point
gain over the $[0,1]$ linear baseline. The quantum kernel achieves a
significant margin over a classical SVM ($\chi^{2}=10.32$, $p=0.0013$),
though its 2.68-point margin over Random Forest is not statistically
significant. The classifier attains complete Myocardial Infarction
separation (AUC\,=\,1.000), Expected Calibration Error below 0.07, and
retains 86.7\% accuracy under brightness perturbation where the
classical SVM falls to 46.7\%. Zero-shot transfer to a second scanner
domain reverses the ordering, indicating that generalisation is
currently bounded by classical feature extraction. The pipeline is
deployed as an end-to-end clinical screening web application.
\end{abstract}

\begin{IEEEkeywords}
Quantum machine learning, quantum kernel methods, ZZFeatureMap, ECG
classification, circuit ablation, statistical validation.
\end{IEEEkeywords}

%-----------------------------------------------------------------------
\section{Introduction}
%-----------------------------------------------------------------------

Cardiovascular diseases are the leading cause of death worldwide,
responsible for approximately 17.9~million deaths annually according to
World Health Organization estimates~\cite{who2021}. The term covers
conditions including myocardial infarction, arrhythmia, coronary artery
disease and stroke. The Electrocardiogram (ECG) remains the standard
first-line, non-invasive tool for detecting these conditions. Specialist
interpretation of ECG traces is frequently unavailable in high-volume
hospitals and rural clinics, which gives an automated classification
system real clinical value.

Machine learning applied to ECG classification has progressed
considerably over the past decade, moving from hand-crafted signal
features toward representations learned directly by deep networks.
Progress specifically on ECG \emph{image} classification has been
slower, because clinical image datasets remain small and multi-class,
and non-linear correlations between waveform regions are difficult for
a compact classical feature space to separate cleanly.

Quantum machine learning offers one route around that limitation. A
quantum feature map encodes classical data into the state of a
multi-qubit system occupying a Hilbert space whose dimension grows
exponentially with qubit count, giving the resulting quantum kernel
access to a far richer set of non-linear decision boundaries.
Present-day hardware remains in the Noisy Intermediate-Scale Quantum
(NISQ) era, so near-term research of this kind is evaluated on
classical simulators rather than physical devices, as is the case
throughout this study.

Prior work applying quantum machine learning to ECG classification has
established that the approach is workable, but it has generally done so
by reporting a single accuracy figure under one fixed circuit
configuration, with an unexamined encoding range, no statistical
validation of the reported margin and no calibration analysis. This
paper addresses that gap directly. The classical portion of the pipeline
is held constant while five circuit design variables are systematically
varied: encoding range, entanglement topology, Pauli operator structure,
circuit repetition depth and the number of retained classical features.

The main contributions of this study are as follows.
\begin{enumerate}
  \item \textbf{Controlled $2\!\times\!2$ factorial circuit ablation.}
    Scaling features to $[0,\pi]$ and closing the entanglement chain into
    a ring each raise accuracy at both levels of the other factor; their
    joint effect is sub-additive by 4.84 against 5.92 percentage points.
    The result is consistent with full $2\pi$ phase coverage and ring
    connectivity acting on an overlapping set of misclassified samples.
  \item \textbf{Statistical validation against classical baselines.}
    McNemar's test is applied to every pairwise model comparison,
    including the strongest classical baseline. To the best of our
    knowledge this is the first significance test of quantum kernel
    advantage reported for ECG image classification.
  \item \textbf{Calibration and perturbation robustness.} Expected
    Calibration Error below 0.07 for all three deployed models, and a
    nine-condition perturbation study in which phase encoding retains
    86.7\% under brightness distortion against 46.7\% for the classical
    SVM.
  \item \textbf{Domain boundary.} A zero-shot cross-scanner evaluation
    identifies feature-extractor domain adaptation as the blocking step
    for deployment beyond the training scanner.
  \item \textbf{Deployed system.} A web application with automated input
    gatekeeping, three-model consensus comparison and PDF report
    generation, together with the code and scripts that reproduce every
    table and figure.
\end{enumerate}

%-----------------------------------------------------------------------
\section{Related Work}
%-----------------------------------------------------------------------

\subsection{Classical ECG Classification}

Yildirim et al.~\cite{yildirim2018} showed that a one-dimensional CNN
applied directly to raw ECG waveforms achieves 91.33\% accuracy across
17 arrhythmia classes from the MIT-BIH database. Vasquez-Iturralde
et al.~\cite{vasquez2024} reviewed 494 deep learning ECG publications
and identified the scarcity of studies operating on image-format ECG
inputs and the near-absence of confidence calibration reporting as
persistent gaps. Working directly with image-format ECGs is distinctly
more challenging: the lack of a digital waveform precludes standard
frequency-domain techniques, requiring classifiers to handle
binarisation artefacts, grid lines and scan-quality variations.
Prabhu et al.~\cite{prabhu2023} report a classical ResNet50--SVM
baseline reaching 83.33\% on the same four-class clinical dataset used
here. The present work provides both calibration and statistical
significance analyses, directly addressing the gaps identified in the
broader literature.

\subsection{Theoretical Foundations of Quantum Kernel Methods}

The theoretical basis for quantum kernel classification was established
by Havlicek et al.~\cite{havlicek2019}, who constructed
ZZFeatureMap-based kernels that map data into exponentially large
Hilbert spaces through Pauli-ZZ interactions. Huang et al.~\cite{huang2021}
showed that once training data is available a classical model can often
match a quantum model, and proposed a geometric criterion for judging
whether a quantum advantage is plausible for a given dataset. Biamonte
et al.~\cite{biamonte2017} identified hybrid classical-quantum pipelines
as the most practical near-term architecture, directly motivating the
Truncated SVD stage described in Section~\ref{sec:method}.

\subsection{Quantum Machine Learning for Cardiac and ECG Classification}

Prabhu et al.~\cite{prabhu2023} demonstrated quantum ECG classification
on the same four-class clinical data used here, combining ResNet50
\texttt{pool1\_pool} features with Truncated SVD at nine components and
a ZZFeatureMap QSVC under $[0,1]$ feature scaling, reporting 94.09\%
for the QSVC. The paper does not state the entanglement topology;
Qiskit's ZZFeatureMap default is full entanglement, so the configuration
used was almost certainly $[0,1]$\,+\,full. Re-running that
configuration on the present pipeline gives 93.01\%, reducing the
reproduction gap to 1.08 points attributable to split and seed
differences. Neither the encoding range nor the entanglement topology
was varied, and no statistical validation, calibration analysis or
deployment was described. The present work retains that architecture and
adds the circuit design study.

Jain et al.~\cite{jain2025} applied ResNet50 with PCA at eight
components and an eight-qubit ZZFeatureMap QSVC to two ECG image
datasets, reporting that the QSVC outperformed a classical SVM by
roughly eight percentage points, with significance assessed by a
non-parametric rank test. Their second dataset is the one used here for
the cross-dataset transfer study in Section~\ref{sec:transfer}. Neither
the encoding range nor the entanglement topology was varied.
Ramkhelawan et al.~\cite{ramkhelawan2025} combined ResNet50 with a
ten-qubit variational quantum circuit for binary arrhythmia detection,
reporting 99.98\% accuracy on 123,998 images; the binary setup is
structurally simpler and the dataset is two orders of magnitude larger.
Ozpolat and Karabatak~\cite{ozpolat2023} found that accuracy does not
increase monotonically with qubit count, consistent with the
circuit-repetition ablation in Section~\ref{sec:ablation}. Aksoy and
Ozpolat~\cite{aksoy2025} adopted $[0,\pi]$ scaling independently on
tabular medical data, supporting reading the $[0,\pi]$ result as a
property of the ZZFeatureMap rotation structure rather than a
dataset-specific artefact. Soon et al.~\cite{soon2026} proposed a
hierarchical quantum ensemble reaching 97\% on tabular cardiovascular
data; the architecture does not involve ECG images and does not ablate
circuit design. Setiawan et al.~\cite{setiawan2026} reviewed 94
quantum machine learning healthcare studies and identified calibration
analysis and statistical significance testing as almost entirely absent
from the literature.

\subsection{Summary and Positioning}

Table~\ref{tab:related} positions the present work against prior
quantum ECG and cardiac studies. Among all works surveyed, none combines
a systematic circuit design ablation, a significance test, a calibration
analysis and a deployed interface on an ECG image classification task.

\begin{table}[t]
\caption{Comparison with Prior Quantum ECG and Cardiac Works}
\label{tab:related}
\centering
\renewcommand{\arraystretch}{1.1}
\resizebox{\columnwidth}{!}{%
\begin{tabular}{lcccccc}
\toprule
Reference & Task & Best Acc. & Ablation & Stat. & Calib. & Depl.\\
\midrule
Yildirim et al.~\cite{yildirim2018} & Arr.\ sig. & 91.33\% & No & No & No & No\\
Havlicek et al.~\cite{havlicek2019} & Synthetic & N/A & No & No & No & No\\
Huang et al.~\cite{huang2021} & Geom.\ anal. & N/A & No & No & No & No\\
Prabhu et al.~\cite{prabhu2023} & 4-cls ECG & 94.09\% & No & No & No & No\\
Jain et al.~\cite{jain2025} & ECG img. & ---$^{\dagger}$ & No & Rank & No & No\\
Ramkhelawan et al.~\cite{ramkhelawan2025} & Bin.\ arr. & 99.98\% & No & No & No & No\\
Ozpolat \& Kar.~\cite{ozpolat2023} & Arrhythmia & 80.90\% & Qubit & No & No & No\\
Aksoy \& Ozpolat~\cite{aksoy2025} & Med.\ tab. & QSVM best & No & No & No & No\\
Soon et al.~\cite{soon2026} & Tab.\ CVD & 97.00\% & No & No & No & No\\
Setiawan et al.~\cite{setiawan2026} & QML rev. & N/A & N/A & N/A & N/A & N/A\\
\textbf{This work} & \textbf{4-cls ECG} & \textbf{94.62\%} & \textbf{320 cfg} & \textbf{McNemar} & \textbf{ECE$<$0.07} & \textbf{Yes}\\
\bottomrule
\end{tabular}}
\smallskip

\begin{flushleft}
{\footnotesize $^{\dagger}$Results reported across two datasets with
different splits; no single comparable figure available.}
\end{flushleft}
\end{table}

%-----------------------------------------------------------------------
\section{Proposed Methodology}
\label{sec:method}
%-----------------------------------------------------------------------

The proposed system is divided into two sequential stages: a Gating
Layer that filters and validates incoming uploads, and a Classification
Engine that extracts features and computes the diagnostic prediction.
Fig.~\ref{fig:pipeline} shows the complete inference data flow. During
development, an offline training phase fits the feature reducers and
trains the classifiers on 928 images, serialising the resulting models
to disk. At inference time, those objects are loaded to evaluate single
uploads without retraining.

\begin{figure}[t]
  \centering
  \fbox{\parbox{0.9\columnwidth}{\centering\vspace{1.5cm}%
    [Pipeline block diagram -- draw.io figure to be inserted here]%
    \vspace{1.5cm}}}
  \caption{System pipeline. The Gating Layer (left) applies three
    sequential validation stages before the Classification Engine
    (right) extracts ResNet50 features, compresses them to 9-D via
    TruncatedSVD, and classifies with three parallel models. Quantum
    branches scale features to $[0,\pi]$ before ZZFeatureMap encoding;
    the classical SVM uses $[0,1]$ features directly.}
  \label{fig:pipeline}
\end{figure}

\subsection{ECG Image Preprocessing}

A seven-step OpenCV pipeline converts each image to a
$340\!\times\!340$ pixel normalised greyscale array. The image is first
converted to greyscale and binarised with Otsu adaptive thresholding.
Morphological dilation with a $15\!\times\!5$ kernel joins fragmented
trace segments, after which the largest contour is detected and the
image is cropped to its bounding box with 15 pixels of padding. A
second Otsu threshold with conditional inversion produces a consistent
trace-on-background polarity across all scan types. Vertical grid lines
are removed with a $1\!\times\!40$ morphological opening and horizontal
grid lines with a $40\!\times\!1$ kernel; the ECG trace is thicker than
any grid line along both axes and therefore survives both subtractions
intact. A $3\!\times\!3$ Gaussian blur and a re-threshold at pixel
value~30 suppress residual speckle, after which the image is resized to
$340\!\times\!340$ using \texttt{cv2.INTER\_AREA} and divided by~255.0
to yield a float32 array in $[0,1]$.

\subsection{ResNet50 Feature Extraction}

Feature extraction uses a ResNet50~\cite{he2016} with frozen ImageNet
weights, rebuilt through the Keras functional API to output at the
\texttt{pool1\_pool} layer. That layer applies $3\!\times\!3$
max-pooling over the initial $7\!\times\!7$ convolutional stage (64
filters, stride~2), producing an $85\!\times\!85\!\times\!64$ spatial
activation map, which is 462,400 dimensions when flattened. The shallow
layer captures low-level morphological structure such as QRS complex
geometry, ST segment curvature and P and T wave shape, without the
high-level ImageNet object semantics encoded at deeper layers. The
choice of extraction layer follows Prabhu et al.~\cite{prabhu2023}.
Greyscale inputs are expanded to three channels by axis repetition
before ResNet50's per-channel normalisation. Extraction ran in batches
of eight.

\subsection{Dimensionality Reduction}

TruncatedSVD at \texttt{n\_components}\,$=$\,9, fitted on the 742
training samples, compresses the 462,400-dimensional vectors to nine
dimensions. TruncatedSVD is used in preference to PCA because it does
not require mean-centering of the data matrix, which is expensive at
this dimensionality and destroys the sparsity of the activation map;
the retained components span directions of largest second moment rather
than of largest variance. The choice of nine components follows Prabhu
et al.~\cite{prabhu2023} and is confirmed by the ablation in
Section~\ref{sec:svd-reps}. A MinMaxScaler, also fitted on the training
data alone, normalises the nine-dimensional vectors to $[0,1]$. For the
quantum classifiers the scaled vector is then multiplied by $\pi$ before
encoding, placing features in $[0,\pi]$; the classical SVM receives the
$[0,1]$ vector directly. Both fitted objects are serialised at training
time and loaded once at API startup.

\subsection{Classical Support Vector Machine}

The classical SVM uses \texttt{sklearn.svm.SVC}~\cite{sklearn} with an
RBF kernel, $C\!=\!10$, \texttt{gamma}\,$=$\,\texttt{`scale'} and
\texttt{class\_weight}\,$=$\,\texttt{`balanced'}. A
\texttt{CalibratedClassifierCV} wrapper with
\texttt{method}\,$=$\,\texttt{`isotonic'} and \texttt{cv}\,$=$\,3
produces calibrated class probabilities. Hyperparameters were selected
by five-fold stratified GridSearchCV maximising macro F1 over
$C\in\{0.01,0.1,1.0,5.0,10.0\}$ and
$\texttt{gamma}\in\{\texttt{`scale'},\texttt{`auto'}\}$. The selected
$C$ lies at the upper edge of the grid, indicating the search was
bounded from above by the grid rather than by the data.

\subsection{Quantum Support Vector Classifier}

The QSVC encodes the nine-dimensional vector into a
$2^9\!=\!512$-dimensional complex Hilbert space using Qiskit's
ZZFeatureMap~\cite{havlicek2019}, the Pauli-ZZ construction of
Havlicek~et~al. Each repetition applies a Hadamard layer across all nine
qubits, a first-order phase gate $P(2x_i)$ on qubit~$i$, and a
second-order block
$\mathrm{CNOT}$--$P\!\left(2(\pi-x_i)(\pi-x_j)\right)$--$\mathrm{CNOT}$
on each connected qubit pair under the selected topology. The fidelity
quantum kernel is
\begin{equation}
K(\mathbf{x},\mathbf{z}) = |\langle\psi(\mathbf{x})|\psi(\mathbf{z})\rangle|^2
\label{eq:kernel}
\end{equation}
where $|\psi(\mathbf{x})\rangle$ is the statevector produced by
encoding input~$\mathbf{x}$. Rather than running $N^2$ circuit
evaluations, all $N\!=\!742$ training statevectors are precomputed once
and the kernel matrix is assembled as
\begin{equation}
K[i,j] = |\mathbf{sv}_i \cdot \overline{\mathbf{sv}_j}|^2
\label{eq:kmatrix}
\end{equation}
where $\mathbf{sv}_i$ is the 512-dimensional complex amplitude vector
for sample~$i$. Full matrix construction for 742 samples takes under
15~seconds on CPU, against over 30~minutes for direct circuit
evaluation. A grid search over four encoding ranges, five entanglement
topologies, $\texttt{reps}\in\{1,2,3,4\}$ and
$C\in\{0.1,1.0,5.0,10.0\}$ (320 configurations) identified $[0,\pi]$
encoding, circular entanglement, $\texttt{reps}\!=\!2$ and $C\!=\!5.0$
as the optimal setting.

\subsection{Multiclass Pegasos QSVC}

Six one-versus-one binary Pegasos QSVC models~\cite{pegasos} cover all
four class pairs. The native Qiskit PegasosQSVC does not converge at
nine qubits: the average kernel value at this scale is approximately
$1/512$, so the accumulated decision function remains numerically close
to zero over the default iteration budget. A custom training loop
precomputes all statevectors and evaluates each gradient step through
matrix multiplication, and a per-model grid search over the
regularisation constant and the iteration count recovers accuracy from
78.49\% to 91.94\%.

\subsection{Gating Layer and Web Application}

The Gating Layer protects the classification pipeline using three
sequential guards. Gate~1 is a MobileNetV2~\cite{sandler2018} model
that hard-rejects non-ECG uploads (HTTP~422). Gate~2 is a Periodicity
Filter that computes the autocorrelation of per-column dark pixel
density, rejecting images with a peak autocorrelation below~0.15 or an
axis-angle gradient fraction above~0.70 (HTTP~422). Gate~3 is a
Mahalanobis distance guard applied after SVD compression; it compares
the 9-D feature vector against the training distribution and appends a
non-blocking out-of-distribution flag to the response if the distance
exceeds the 99th-percentile threshold of~4.59.

The FastAPI backend exposes six REST endpoints. A Redis cache keyed by
the SHA-256 hash of the image bytes together with the model identifier
returns stored results for repeated submissions with a 24-hour
time-to-live. The React frontend provides tabs for Diagnosis, Batch
Upload, Model Performance, Quantum Insights, API documentation and
prediction history, with colour-coded severity indicators and calibrated
confidence bars. Measured end-to-end QSVC latency is 893.9~ms
(preprocessing 126.9~ms, ResNet50 627.6~ms, SVD 52.0~ms,
classification 87.4~ms).

%-----------------------------------------------------------------------
\section{Experimental Setup}
%-----------------------------------------------------------------------

The experiments use the ECG image collection released on Mendeley Data
by Khan et al.~\cite{khan2021}, collected at the Ch.\ Pervaiz Elahi
Institute of Cardiology, Multan, Pakistan. The standard four-class
subset comprises 929~images (Normal~284, Arrhythmia~233, Myocardial
Infarction~240, History of MI~172), consistent with~\cite{prabhu2023}.
One Myocardial Infarction scan could not be decoded and was excluded,
leaving 928~images. Table~\ref{tab:dataset} shows the class
distribution and the stratified 80:20 split applied with
\texttt{random\_state}\,$=$\,42. No augmentation was applied during
training; augmented images appear only in the robustness evaluation of
Section~\ref{sec:robust}.

\begin{table}[t]
\caption{Dataset Class Distribution and Train/Test Split}
\label{tab:dataset}
\centering
\renewcommand{\arraystretch}{1.1}
\begin{tabular}{lrrr}
\toprule
Class & Total & Train & Test\\
\midrule
Normal              & 284 & 227 & 57\\
Arrhythmia          & 233 & 186 & 47\\
Myocardial Infarction & 239 & 191 & 48\\
History of MI       & 172 & 138 & 34\\
\midrule
\textbf{Total}      & \textbf{928} & \textbf{742} & \textbf{186}\\
\bottomrule
\end{tabular}
\end{table}

All models were evaluated on the same held-out 186-image test set. The
TruncatedSVD reducer and the MinMaxScaler were fitted on the 742
training samples and serialised; the test set passed through the same
fitted objects without refitting, preventing the leakage failure mode
identified in~\cite{vasquez2024}. The primary metrics are accuracy and
macro-averaged F1, with per-class and macro-averaged ROC-AUC also
reported. McNemar's test~\cite{mcnemar1947} is applied to all pairwise
model comparisons. Expected Calibration Error (ECE) is computed with
15 equal-width probability bins.

Because the test set contains 186 images, a single reclassified sample
shifts accuracy by 0.54 percentage points. Differences of this
magnitude are reported for completeness but are not interpreted as
evidence of a real performance gap. The ablation effects discussed in
Section~\ref{sec:ablation} range from 4 to 26 percentage points,
corresponding to 8 to 49 samples.

All experiments ran on a standard CPU (Intel Core~i7, 16~GB RAM).
Quantum circuits were simulated with the Qiskit Aer~0.13 statevector
simulator~\cite{qiskit2023}. The software stack uses Python~3.10,
TensorFlow~2.15, Scikit-learn~1.3~\cite{sklearn} and Qiskit Machine
Learning~0.7. Random Forest and k-Nearest Neighbour classifiers were
also trained on the same nine-dimensional MinMax-scaled features and
are included in Table~\ref{tab:perf} as additional classical reference
points.

Kernel Target Alignment (KTA)~\cite{kta2001} was computed for every
distinct kernel in the ablation sweep. Because KTA is a property of the
kernel matrix and the labels alone, it does not vary with the SVM
regularisation constant, so the 320-configuration sweep contains
64~distinct kernels across four encoding ranges, four entanglement
topologies and four repetition counts. KTA values span 0.1063 to
0.1277. The highest value, 0.1277, is attained by $[0,1]$ encoding with
full entanglement at $\texttt{reps}\!=\!1$, which is not the
highest-accuracy configuration. The selected circuit, $[0,\pi]$ with
circular entanglement at $\texttt{reps}\!=\!2$, has KTA\,=\,0.1131,
mid-range and 0.0146 below the maximum. The two highest-KTA
configurations, both raw $[0,1]$ with full entanglement at
$\texttt{reps}\!=\!1$, produce 89.78\% and 92.47\% respectively, well
below the 94.62\% of the selected circuit. Across the sweep the KTA
ranking and the accuracy ranking disagree on both the best encoding
range and the best topology, so higher alignment does not imply higher
accuracy in this setting. That is a useful negative result for
practitioners, since KTA is frequently proposed as a cheap
circuit-selection criterion that avoids training a classifier per
configuration. It is reported here as a descriptive statistic for the
selected circuit rather than as a selection rule.

%-----------------------------------------------------------------------
\section{Results and Discussion}
%-----------------------------------------------------------------------

\subsection{Model Performance}

Table~\ref{tab:perf} reports accuracy, precision, recall and macro F1
for all five classifiers on the held-out 186-image test set.

\begin{table}[t]
\caption{Model Performance on the 186-Image Test Set}
\label{tab:perf}
\centering
\renewcommand{\arraystretch}{1.1}
\resizebox{\columnwidth}{!}{%
\begin{tabular}{lccccr}
\toprule
Model & Acc. & Prec. & Recall & F1 & \cite{prabhu2023}\\
\midrule
Classical SVM & 84.95\% & 84.17\% & 84.41\% & 84.10\% & 83.33\%\\
\textbf{QSVC (ours)} & \textbf{94.62\%} & \textbf{94.84\%} & \textbf{94.05\%} & \textbf{94.39\%} & 94.09\%\\
Pegasos QSVC & 91.94\% & 91.99\% & 91.61\% & 91.31\% & 93.05\%\\
Random Forest & 91.94\% & 91.39\% & 91.80\% & 91.20\% & N/A\\
k-NN & 84.95\% & 85.54\% & 85.07\% & 83.61\% & N/A\\
\bottomrule
\end{tabular}}
\end{table}

The tuned QSVC leads the four other classifiers. Its margin over the
classical SVM is 9.67 percentage points, but the SVM is not the
strongest classical baseline in the table. Random Forest reaches 91.94\%
on the same nine-dimensional features with no quantum component, which
reduces the margin that the quantum kernel actually has to defend to
2.68 points, or five test images. First, the nine-dimensional SVD
features already carry strong discriminative signal, so the quantum
kernel is improving on an informative representation rather than
compensating for a weak one. Second, any claim of quantum advantage must
be stated against Random Forest as well as against the SVM, which
Section~\ref{sec:sig} does.

The classical SVM exceeds the corresponding baseline in~\cite{prabhu2023}
by 1.62 points, which we tentatively attribute to balanced class
weighting and to the resampling performed by the calibration wrapper.
Pegasos falls 1.11 points below the 93.05\% reported
in~\cite{prabhu2023}; the likely cause is stricter leakage control here.
The 0.53-point margin over the QSVC figure reported in~\cite{prabhu2023}
corresponds to a single test image, and no conclusion is drawn from it.

Table~\ref{tab:perclass} gives per-class precision, recall and F1 for
the three deployed classifiers.

\begin{table}[t]
\caption{Per-Class Precision / Recall / F1 (186-Image Test Set)}
\label{tab:perclass}
\centering
\renewcommand{\arraystretch}{1.1}
\resizebox{\columnwidth}{!}{%
\begin{tabular}{lcccc}
\toprule
Class & SVM P/R/F1 & QSVC P/R/F1 & Pegasos P/R/F1 & $n$\\
\midrule
Normal    & 0.870/0.825/0.847 & \textbf{0.948/0.965/0.957} & 0.891/\textbf{1.000}/0.942 & 57\\
Arrhythmia & 0.881/0.787/0.832 & \textbf{0.878/0.915/0.896} & 0.944/0.723/0.819 & 47\\
MI        & 0.873/\textbf{1.000}/0.932 & \textbf{1.000/1.000/1.000} & 0.980/\textbf{1.000}/0.990 & 48\\
Hist.\ MI & 0.743/0.765/0.754 & \textbf{0.968/0.882/0.923} & 0.865/0.941/0.901 & 34\\
\bottomrule
\end{tabular}}
\end{table}

The QSVC achieves 100\% recall on Myocardial Infarction (48 of 48), the
class for which a missed detection carries the greatest clinical cost;
all three classifiers reach 100\% recall on this class. The QSVC leads
on every class except Normal recall, where Pegasos achieves perfect
recall at the cost of some precision. Arrhythmia is the most frequently
confused class across all models, consistent with the morphological
variability of rhythm abnormalities. History of MI shows the largest
gap between the QSVC (F1\,=\,0.923) and the classical SVM
(F1\,=\,0.754), indicating that the quantum kernel's advantage is
concentrated on the two most ambiguous classes.

\subsection{Statistical Significance}
\label{sec:sig}

Table~\ref{tab:mcnemar} reports McNemar's test for all pairwise model
comparisons.

\begin{table}[t]
\caption{McNemar's Test Results ($n=186$ Test Samples)}
\label{tab:mcnemar}
\centering
\renewcommand{\arraystretch}{1.1}
\resizebox{\columnwidth}{!}{%
\begin{tabular}{lccccl}
\toprule
Comparison & $b$ & $c$ & $\chi^2$ & $p$ & Sig.\\
\midrule
\textbf{QSVC vs.\ SVM} & \textbf{23} & \textbf{5} & \textbf{10.321} & \textbf{0.0013} & \textbf{Yes}\\
Pegasos vs.\ SVM & 19 & 6 & 5.760 & 0.0164 & Yes\\
QSVC vs.\ Pegasos & 9 & 4 & 1.231 & 0.267 & No\\
QSVC vs.\ Rand.\ Forest & 11 & 6 & 0.941 & 0.332 & No\\
\bottomrule
\end{tabular}}
\end{table}

The QSVC corrected 23 samples that the classical SVM misclassified
while missing only~5 that the SVM classified correctly. With continuity
correction, $\chi^2=(|23-5|-1)^2/(23+5)=10.32$ at $p=0.0013$, which
indicates that the difference is a systematic property of the two
models rather than an artefact of the particular split. Both quantum
models are significantly better than the classical SVM at the 0.05
threshold.

The remaining rows bound the claim. The QSVC versus Random Forest
comparison gives $\chi^2=0.941$ at $p=0.332$, well short of
significance. The quantum kernel is significantly better than a tuned
RBF SVM on this data, though better than the strongest classical
baseline by a margin this test set cannot resolve. To the best of our
knowledge, no prior study of ECG image classification reports a
significance test of quantum kernel advantage at all.

\subsection{ROC-AUC and Confidence Calibration}
\label{sec:roc-roc}

Both quantum models reach AUC\,=\,1.000 on Myocardial Infarction,
ranking every MI sample above every non-MI sample at every threshold.
The QSVC macro ROC-AUC of 0.989 exceeds the classical SVM's 0.940 by
0.049 across all four classes. Expected Calibration Error is below 0.07
for all three models: 0.054 for the SVM, 0.063 for the QSVC and 0.067
for Pegasos. A model reporting 80\% confidence is therefore correct in
approximately 80\% of cases, a property required before a system can
meaningfully support a clinical decision. Calibration analysis of this
kind was not found in any of the ECG quantum machine learning papers
surveyed in Section~II.

\subsection{Augmentation Robustness}
\label{sec:robust}

\begin{figure}[t]
  \centering
  \fbox{\parbox{0.9\columnwidth}{\centering\vspace{1.5cm}%
    [Figure: augmentation robustness bar chart -- to be inserted here]%
    \vspace{1.5cm}}}
  \caption{Model accuracy under nine augmentation conditions. The
    classical SVM falls to 46.7\% under brightness perturbation while
    the QSVC maintains 86.7\%.}
  \label{fig:augmentation}
\end{figure}

The robustness evaluation used a fixed subset of 45~images from a
second ECG collection, covering Normal and Arrhythmia cases only, held
constant across every perturbation condition. The subset spans only two
of the four training classes, and it is small enough that a single
image moves accuracy by 2.2~points. Unperturbed, both the classical SVM
and the QSVC classify 38 of 45 (84.4\%). Under brightness perturbation,
which simulates an overexposed photograph of a paper ECG print, the
classical SVM falls to 21~of~45 (46.7\%) while the QSVC classifies
39~of~45 (86.7\%). The 40.0-point gap is far larger than the 2.2-point
resolution of a single image, which is why the direction of the effect
is reported; its magnitude is not precisely estimated.

The ZZFeatureMap encodes feature values as phase angles, and phase is
periodic, so a uniform brightness shift advances the encoded state
around the Bloch sphere rather than displacing it away from the training
support vectors. The RBF kernel operates on Euclidean distance in the
same nine-dimensional compressed space, where the same shift translates
every test sample away from its nearest training neighbours. Brightness
is the one perturbation axis where phase encoding produces a measurable
advantage, and it is also the axis that varies most in practice, since
ward ECG prints are photographed under whatever light is available.

\subsection{Cross-Dataset Transfer}
\label{sec:transfer}

The three trained models were applied without retraining to the second
Mendeley ECG collection used by Jain et al.~\cite{jain2025}, comprising
707~images across Normal~(295), Abnormal Heartbeat/Arrhythmia~(241) and
History of MI~(171). The Myocardial Infarction class is absent from
that collection, so any prediction of it counts as an error.

\begin{table}[t]
\caption{Zero-Shot Transfer to a Second Scanner Domain
         (707~Images, Three Classes)}
\label{tab:transfer}
\centering
\renewcommand{\arraystretch}{1.1}
\begin{tabular}{lccc}
\toprule
Model & 3-cls Acc. & 3-cls F1 & Binary Acc.\\
\midrule
Classical SVM        & \textbf{62.66\%} & \textbf{0.638} & \textbf{70.72\%}\\
QSVC (frozen)        & 34.09\% & 0.173 & 58.27\%\\
Pegasos QSVC (frozen) & 45.40\% & 0.435 & 60.25\%\\
\bottomrule
\end{tabular}
\end{table}

The ordering reverses. The frozen ResNet50 was never adapted to either
scanner, and the second collection differs in contrast, brightness
distribution and trace density, which relocates the SVD-compressed
vectors into a region the training support vectors cover poorly. The
consequence is specific: the 94.62\% figure, the McNemar result and the
brightness robustness are all properties of the training scanner domain,
and domain adaptation of the feature extractor, not a change to the
quantum circuit, is the step that would extend them.

%-----------------------------------------------------------------------
\section{Ablation Studies}
\label{sec:ablation}
%-----------------------------------------------------------------------

\subsection{Feature Encoding Range}

Table~\ref{tab:encoding} compares QSVC and Pegasos accuracy across four
encoding ranges. All rows use circular entanglement,
$\texttt{reps}\!=\!2$ and $C\!=\!5.0$.

\begin{table}[t]
\caption{Accuracy vs.\ Feature Encoding Range\\
         (all rows: circular entanglement, reps\,=\,2, $C$\,=\,5.0)}
\label{tab:encoding}
\centering
\renewcommand{\arraystretch}{1.1}
\begin{tabular}{lrr}
\toprule
Encoding & QSVC Acc. & Pegasos Acc.\\
\midrule
L2 normalisation (unit sphere) & 68.28\% & 57.53\%\\
MinMax $[0,1]$                 & 92.47\% & 84.41\%\\
\textbf{MinMax $[0,\pi]$ (ours)} & \textbf{94.62\%} & \textbf{90.86\%}$^{\dagger}$\\
MinMax $[0,2\pi]$              & 94.09\% & 90.32\%\\
\bottomrule
\end{tabular}
\smallskip

\begin{flushleft}
{\footnotesize $^{\dagger}$Pegasos at shared $C\!=\!5.0$. The final
tuned result (91.94\%, Table~\ref{tab:perf}) uses per-model
regularisation and iteration tuning.}
\end{flushleft}
\end{table}

\begin{figure}[t]
  \centering
  \fbox{\parbox{0.9\columnwidth}{\centering\vspace{1.5cm}%
    [Figure: QSVC accuracy vs.\ encoding range bar chart -- to be
    inserted here]%
    \vspace{1.5cm}}}
  \caption{QSVC accuracy across four feature encoding ranges, all with
    circular entanglement. L2 normalisation discards feature magnitude
    and produces the largest drop. Moving to $[0,2\pi]$ introduces
    phase wrap-around.}
  \label{fig:encoding}
\end{figure}

After the Hadamard layer the qubit states lie on the equator of the
Bloch sphere, and the $R_z$ gates advance the azimuthal phase angle. A
feature confined to $[0,1]$ sweeps approximately~2~radians of the
$2\pi$ available, so the encoded states remain clustered within a narrow
arc. Extending the range to $[0,\pi]$ opens the sweep to the full
circle, distributing the encoded states more widely. This moves away
from the kernel concentration regime analysed by Thanasilp
et al.~\cite{thanasilp2024}, in which kernel values collapse towards a
fixed point and the trained model becomes insensitive to its input.
Moving from $[0,1]$ to $[0,\pi]$ raises the effective rank of the
kernel matrix by~38\% (from 94.6 to 130.4) and accuracy by 2.15
percentage points under circular entanglement. L2 normalisation performs
worst by a wide margin (26.34~points below $[0,\pi]$), because
projecting onto the unit sphere discards feature magnitude entirely, and
magnitude is what the dominant uncentred SVD components encode.

To check whether a range marginally below~$\pi$ would avoid the
endpoint-collision concern and equal or exceed~$\pi$, a fine sweep at
six sub-$\pi$ values was run under circular entanglement
($\texttt{reps}\!=\!2$, $C\!=\!5.0$).

\begin{table}[t]
\caption{Encoding Range Fine Sweep\\
         (circular entanglement, reps\,=\,2, $C$\,=\,5.0)}
\label{tab:finesweep}
\centering
\renewcommand{\arraystretch}{1.1}
\begin{tabular}{lrr}
\toprule
Range & Accuracy & Correct\,/\,186\\
\midrule
$[0,\,0.80\pi]$ & 92.47\% & 172\\
$[0,\,0.85\pi]$ & 93.55\% & 174\\
$[0,\,0.90\pi]$ & 92.47\% & 172\\
$[0,\,0.95\pi]$ & 91.40\% & 170\\
$\mathbf{[0,\,\pi]}$ & \textbf{94.62\%} & \textbf{176}\\
$[0,\,1.05\pi]$ & 94.09\% & 175\\
\bottomrule
\end{tabular}
\end{table}

No sub-$\pi$ value exceeds $[0,\pi]$. The best of them, $0.85\pi$ at
174 of~186, is two test images below~$\pi$, and $1.05\pi$ is one image
below. The endpoint collision therefore does not cost measurable
accuracy on this dataset.

The quantity that tracks accuracy here is the effective rank rather than
the dispersion of kernel values. Moving from $[0,1]$ to $[0,\pi]$
raises the effective rank by~38\% alongside the 2.15-point accuracy
gain, so a larger fraction of the 512-dimensional space is used to
separate the samples. The dispersion moves the other way: the standard
deviation of off-diagonal kernel values falls by~30\%, and the mean
sits closer to the $1/2^9\!=\!0.00195$ random-overlap floor. The
widened encoding range does not spread the kernel values out; it
redistributes them into more independent directions.

\subsection{Entanglement Topology}

Table~\ref{tab:topology} compares QSVC accuracy across five entanglement
topologies. All rows use $[0,\pi]$ encoding, $\texttt{reps}\!=\!2$ and
$C\!=\!5.0$.

\begin{table}[t]
\caption{Accuracy vs.\ Entanglement Topology\\
          ($[0,\pi]$ encoding, reps\,=\,2, $C$\,=\,5.0)}
\label{tab:topology}
\centering
\renewcommand{\arraystretch}{1.1}
\begin{tabular}{lrl}
\toprule
Topology & QSVC Acc. & Gain vs.\ Linear\\
\midrule
Linear                        & 93.01\% & N/A\\
Pairwise                      & 93.01\% & 0.00~pp\\
Full                          & 92.47\% & $-$0.54~pp\\
\textbf{Circular (ours)}      & \textbf{94.62\%} & $+$1.61~pp\\
Shifted-circular-alt.\ (SCA)  & 94.62\% & $+$1.61~pp\\
\bottomrule
\end{tabular}
\end{table}

The circular topology closes the qubit chain by connecting qubit~8 back
to qubit~0, adding exactly one entangling pair at negligible additional
depth. It is the only topology that stands clearly apart. Full
entanglement does not improve on linear or pairwise at $[0,\pi]$
encoding, plausibly because the much larger number of entangling gates
over-parameterises the kernel for a nine-dimensional input.

\subsection{Separating the Two Factors}

Table~\ref{tab:factorial} gives the four cells of the complete
$2\!\times\!2$ factorial design ($\texttt{reps}\!=\!2$, $C\!=\!5.0$).

\begin{table}[t]
\caption{QSVC Accuracy by Encoding Range and Topology\\
         (reps\,=\,2, $C$\,=\,5.0)}
\label{tab:factorial}
\centering
\renewcommand{\arraystretch}{1.3}
\begin{tabular}{lcc}
\toprule
 & Linear & Circular\\
\midrule
$[0,1]$ & 89.78\% (167/186) & 92.47\% (172/186)\\
$[0,\pi]$ & 93.01\% (173/186) & \textbf{94.62\% (176/186)}\\
\bottomrule
\end{tabular}
\end{table}

The four simple effects follow directly. Changing the topology from
linear to circular is worth 2.69~points under $[0,1]$ encoding and
1.61~points under $[0,\pi]$. Changing the encoding range from $[0,1]$
to $[0,\pi]$ is worth 3.23~points under linear entanglement and
2.15~points under circular. The two effects are sub-additive: moving
from $[0,1]$+linear to $[0,\pi]$+circular gains 4.84~points, which is
1.08~points less than the sum of the two individual gains (5.92). A
five-seed cross-validation confirmed that the cell ordering holds across
all seeds and that the sub-additivity replicates at the mean (5.81 joint
against 6.46 summed). Seed~42 lies 1.1 to 2.1 standard deviations above
the five-seed mean in every cell, so the absolute accuracies quoted
throughout this paper should be read as the upper end of the attainable
range rather than as expected values.

Kernel Target Alignment (KTA)~\cite{kta2001} was computed for all
64~distinct kernels in the ablation sweep. KTA values span 0.1063
to~0.1277; the highest-KTA configurations are not the
highest-accuracy ones. KTA therefore does not serve as a reliable
circuit-selection criterion in this setting, which is a useful negative
result for practitioners since KTA is frequently proposed as a cheap
circuit-selection heuristic.

\subsection{Pauli Feature Map Comparison}

The ZZFeatureMap outperforms all three alternative Pauli maps under
identical conditions (nine qubits, $\texttt{reps}\!=\!2$, $[0,\pi]$
encoding, $C\!=\!5.0$, circular entanglement). The ZFeatureMap, which
applies only single-qubit $R_z$ rotations with no entanglement, drops
to 87.63\%, 6.99~percentage points below the ZZFeatureMap, confirming
the direct contribution of the ZZ entangling interactions. The
$\mathrm{X}$+$\mathrm{XX}$ variant reaches 91.94\% and the
$\mathrm{Y}$+$\mathrm{YY}$ variant falls to 83.33\%, below the
classical SVM baseline.

\subsection{Retained Components and Circuit Repetitions}
\label{sec:svd-reps}

Sweeping the retained components from 2 to 18 shows accuracy peaking at
nine, confirming the choice made in~\cite{prabhu2023}. Below nine,
discriminative structure from the ResNet50 activation map is discarded;
above nine, additional components contribute more noise than signal
relative to the encoding capacity. Sweeping circuit repetitions
confirms $\texttt{reps}\!=\!2$: $\texttt{reps}\!=\!3$ and
$\texttt{reps}\!=\!4$ give no accuracy gain while adding statevector
computation cost, consistent with the finding of Ozpolat and
Karabatak~\cite{ozpolat2023} that increasing circuit depth beyond a
sufficient encoding level does not improve quantum SVM performance.

%-----------------------------------------------------------------------
\section{Limitations}
%-----------------------------------------------------------------------

\textbf{Single split, no variance estimate.} All results come from one
stratified split at $\texttt{random\_state}\!=\!42$. A five-seed
cross-validation is presented in Section~\ref{sec:ablation} but does
not replace a full bootstrap; repeating the full 320-configuration sweep
across ten or more seeds is the first item of further work.

\textbf{Pipeline-relative accuracies.} Re-running $[0,1]$+full
entanglement, the Prabhu et al.~\cite{prabhu2023} configuration as
inferred from the Qiskit default, gives 93.01\% against the 94.09\%
reported there, a residual gap of 1.08~points well within what split
and seed variation can explain. All ablation gains are measured against
configurations run on this pipeline, keeping circuit design comparisons
internally consistent.

\textbf{Domain-bounded advantage.} The zero-shot transfer in
Section~\ref{sec:transfer} reverses the model ordering, so every
accuracy, significance and robustness result is a within-domain result
for the training scanner.

\textbf{Margin against the strongest baseline.} The significance result
is established against a tuned RBF SVM. Against Random Forest, the
margin is 2.68~points and is not resolvable at $n\!=\!186$.

\textbf{Ideal simulation.} All quantum results come from noiseless
statevector simulation, so the reported accuracies are an upper bound on
what a NISQ device would deliver at nine qubits.

\textbf{Inherited extraction layer.} The \texttt{pool1\_pool} layer was
adopted from~\cite{prabhu2023} rather than selected empirically, and no
layer-wise ablation against deeper ResNet50 outputs was run.

%-----------------------------------------------------------------------
\section{Conclusion}
%-----------------------------------------------------------------------

This work treats the quantum feature map as a design space to be
measured rather than a component to be adopted. A controlled
$2\!\times\!2$ factorial ablation over encoding range and entanglement
topology shows that scaling features to $[0,\pi]$ and closing the qubit
chain into a ring each raise accuracy at both levels of the other
factor, and that their joint gain of 4.84~percentage points falls
1.08~points short of the sum of the separate gains. The best of the
four cells, $[0,\pi]$ with circular entanglement at
$\texttt{reps}\!=\!2$, classifies 94.62\% of a held-out 186-image test
set, and McNemar's test places that result significantly above a tuned
classical RBF SVM ($\chi^2\!=\!10.32$, $p\!=\!0.0013$). Expected
Calibration Error below 0.07 for all three deployed models and a
nine-condition perturbation study in which phase encoding retains 86.7\%
under brightness distortion where the classical SVM reaches 46.7\%
accompany it.

Against Random Forest the margin is 2.68~points and does not reach
significance at this test-set size ($b\!=\!11$, $c\!=\!6$,
$\chi^2\!=\!0.941$, $p\!=\!0.332$). Under zero-shot transfer to a
second scanner domain the ordering reverses entirely. These results
locate the contribution precisely: the circuit design choices measured
here are the right ones for this pipeline and this domain, and the
obstacle to generalisation lies in the classical feature extractor
rather than in the quantum circuit. Three items of work address this
directly: domain-adaptive fine-tuning of the ResNet50 extractor,
noise-model evaluation at nine qubits, and gradient-based saliency
overlays that would show clinicians which region of a trace drove a
decision.

%-----------------------------------------------------------------------
\section*{Code and Data Availability}
%-----------------------------------------------------------------------

The ECG image dataset is publicly available on Mendeley
Data~\cite{khan2021}. The implementation, trained model artefacts and
the scripts reproducing all tables and figures are available at
\url{https://github.com/karthikspoojary/QuCardio}.

%-----------------------------------------------------------------------
\section*{Acknowledgement}
%-----------------------------------------------------------------------

The authors thank the Qiskit open-source community and IBM Quantum for
the quantum computing framework. The ECG dataset was made publicly
available by Khan et al.\ through Mendeley Data, collected at the
Ch.\ Pervaiz Elahi Institute of Cardiology, Multan. The authors
acknowledge the support of the Department of Computer Science and
Engineering, St.\ Joseph Engineering College, Mangaluru.

%-----------------------------------------------------------------------
\bibliographystyle{IEEEtran}
\bibliography{references}

\end{document}
